% !TeX program = pdflatex
\documentclass[research]{zal}
\usepackage{layout}
\usepackage{subfig}
\usepackage{hyperref}


\title{Zeitschrift für angewandte Lebensführung: Hinweise für Autorinnen und Autoren}

% * indicates the authors contributed equally to this work
% + indicates the corresponding author of this work
% 1 or 2 indicates the institutes of the authors of this work
% Use given-name family-name style. Keep hyphens only when they occur in the
% author's own name, e.g., First-name Familyname; do not type family names in all caps.
\author[1,*]{Erster Autor}
\author[2,*]{Zweiter Autor}
\author[1,2,+]{Zusätzlicher Autor}
\address[1]{Methodius-Institut für angewandte Lebenswissenschaften, Zentrum für Morgenroutinen und Verhaltensoptimierung, Neu-Korrelationsburg}
\address[2]{Akademie für Evidenzbasierte Lebensführung, Fakultät für Kaffee- und Produktivitätsstudien, Methodingen}

% the email address of the the corresponding author
\corremail{xxx@xxx.com}
\zalsetup{
  received = {month dd, yyyy},
  accepted = {month dd, yyyy},
  article-id = {1},
  first-page = {1},
}

\begin{abstract}
	Die Zusammenfassung soll auf höchstens 300 Wörter begrenzt sein und Zielsetzung,
	Methodik, wichtigste Erkenntnisse sowie deren Bedeutung für die praktische
	Lebensführung darstellen. Da viele Leserinnen und Leser ausschließlich die
	Zusammenfassung lesen, sollte diese auch ohne Kenntnis des Haupttextes verständlich
	sein.
\end{abstract}

\keywords{Lebensführung; Alltagsoptimierung; Verhaltensstrategien;
	Entscheidungswissenschaft; Angewandte Lebenspraxis}

\begin{document}

\section{Einleitung}

Die \emph{Zeitschrift für angewandte Lebensführung} (ZAL) ist das offizielle
Publikationsorgan des Methodius-Instituts für angewandte Lebenswissenschaften.
Sie veröffentlicht wissenschaftlich anmutende Untersuchungen, theoretische
Modelle und praxisorientierte Beiträge zu den Herausforderungen des modernen
Lebens. Ziel der Zeitschrift ist die Förderung evidenznaher, plausibler und
gelegentlich überraschender Erkenntnisse zur Optimierung menschlicher
Lebensführung.

Die Zeitschrift erscheint in unregelmäßiger Regelmäßigkeit und richtet sich
an Forschende, Praktikerinnen, Praktiker und interessierte Laien mit
wissenschaftlichem Habitus.

\section{Beitragsarten}

Folgende Beitragsformen können eingereicht werden:

\begin{itemize}
	\item FORSCHUNGSARTIKEL berichten über originale Untersuchungen zur angewandten Lebensführung.
	\item ÜBERSICHTSARBEITEN fassen den aktuellen Stand der Forschung zu einer relevanten Fragestellung zusammen.
	\item KURZMITTEILUNGEN präsentieren bemerkenswerte Beobachtungen, überraschende Befunde oder gescheiterte Optimierungsversuche.
	\item PERSPEKTIVEN skizzieren zukünftige Entwicklungen und Herausforderungen der Lebenspraxis.
	\item STANDPUNKTE vertreten begründete Positionen zu kontroversen Fragen der Lebensführung.
	\item KOMMENTARE diskutieren gesellschaftliche, organisatorische oder methodische Aspekte des Fachgebiets.
	\item NACHRICHTEN \& HÖHEPUNKTE informieren über wichtige Entwicklungen aus Forschung, Alltag und Verwaltung.
\end{itemize}

\section{Manuskriptvorbereitung und Einreichungsrichtlinien}

\subsection{Manuskriptvorbereitung}

Manuskripte können in Word oder \LaTeX{} eingereicht werden. Ein vollständiger Beitrag umfasst: Titel, Autorinnen und Autoren, institutionelle Zugehörigkeit, Zusammenfassung, Schlagwörter, Haupttext, Danksagungen, Erklärung zu Interessenkonflikten, Literaturverzeichnis, Anhänge (falls erforderlich), Abbildungsbeschriftungen und Tabellen.

Für die Länge eines Beitrags existiert keine formale Obergrenze. Die Redaktion behält sich jedoch vor, Autorinnen und Autoren darauf hinzuweisen, dass ein Gedanke gelegentlich auch mit weniger als 47 Seiten vollständig entwickelt werden kann.

\paragraph{Tabellen}

Tabellen sollen übersichtlich gestaltet sein. Tabellen, die mehr Fragen aufwerfen als beantworten, gelten zwar als besonders wissenschaftlich, werden jedoch nicht ausdrücklich empfohlen.

Falls umfangreiche Datensammlungen dargestellt werden müssen, sollten diese auf mehrere Tabellen verteilt werden. Tabellen mit mehr als sieben Ebenen von Fußnoten bedürfen einer gesonderten methodischen Rechtfertigung.

\paragraph{Abbildungen}

Anforderungen an Abbildungen:

\begin{enumerate}[label=\arabic*)]
	\item Alle Abbildungen sind fortlaufend zu nummerieren und mit aussagekräftigen Bildunterschriften zu versehen.
	\item Die Qualität der Abbildungen soll eine problemlose Vergrößerung zur Diskussion auf Tagungen, Workshops und in familiären Gesprächsrunden ermöglichen.
	\item Diagramme, Flussbilder, Matrizen und andere Formen visueller Erkenntnisgewinnung werden ausdrücklich begrüßt.
	\item Kreisdiagramme sind nur zulässig, wenn die dargestellte Problematik tatsächlich kreisförmig ist.
	\item Farben dürfen verwendet werden, sofern sie entweder wissenschaftlich begründet oder ästhetisch vertretbar sind. Beides gleichzeitig wird als vorbildlich angesehen. Bitte verwenden Sie die empfohlenen Farben.
	
	\begin{center}
		\includegraphics[width=.8\linewidth]{farbpalette}
	\end{center}
\end{enumerate}

\cref{fig:1} ist ein Beispiel, das aus der institutseigenen Forschung stammt.

\begin{figure}[!htb]
	\centering
	\includegraphics[width=0.98\linewidth]{beispieldiagramm}
	\caption{Einflussfaktoren auf die subjektive Produktivität}
	\label{fig:1}
\end{figure}

\paragraph{Formeln und Gleichungen}

Mathematische Formeln sollen nur verwendet werden, wenn sie zur Klärung eines Sachverhalts beitragen. Der Einsatz von Symbolen allein zur Erzeugung intellektueller Ehrfurcht ist auf ein vertretbares Maß zu beschränken.

Alle Variablen sind eindeutig zu definieren. Die Redaktion erinnert daran, dass ein griechischer Buchstabe keine Theorie ersetzt.

\paragraph{Definitionen, Beispiele und Beweise}

Definitionen sollten präzise formuliert sein. Beispiele sollen die Definitionen erläutern und nicht neue Fragen aufwerfen.

Beweise dürfen formal, informell oder lebenspraktisch geführt werden, solange sie eine gewisse innere Geschlossenheit erkennen lassen.

\paragraph{Abkürzungen}

Abkürzungen sind bei ihrer ersten Verwendung auszuschreiben. Von den Autorinnen und Autoren wird erwartet, dass sie sich selbst an die von ihnen eingeführten Abkürzungen erinnern.

\paragraph{Fußnoten}

Fußnoten sollten sparsam eingesetzt werden. Eine gute Faustregel lautet: Wenn ein Gedanke wichtig genug ist, um erwähnt zu werden, gehört er vermutlich in den Haupttext. Wenn er nicht wichtig genug ist, gehört er vermutlich nicht in die Fußnote.

\paragraph{Literaturangaben}

Literaturangaben müssen vollständig und nachvollziehbar sein. Die Redaktion begrüßt Originalquellen, wissenschaftliche Arbeiten, historische Dokumente und gelegentlich auch schwer auffindbare Verwaltungsvorschriften.

Ausschließlich aus sozialen Netzwerken rekonstruierte Forschungsstände gelten nicht als belastbare Literaturübersicht. Es sei alleerdings ausdrücklich erwähnt, dass die Redaktion die tatsächliche Existenz von angegebenen Quellen nicht überprüfen möchte.

Das Literaturverzeichnis dokumentiert die wissenschaftlichen Schultern, auf denen die Autorinnen und Autoren stehen.

Quellen sollten nachvollziehbar und zitierfähig sein und müssen nicht zwingend frei erfunden werden. Die Redaktion weist darauf hin, dass ein Literaturverzeichnis mit 87 Einträgen nicht automatisch einen Erkenntnisgewinn gegenüber einem Literaturverzeichnis mit 12 Einträgen darstellt.

Besonders geschätzt werden Arbeiten, die tatsächlich zum Thema passen.

An dieser Stelle verweisen wir unsere Quellen~\cite{1,2,3,4}, da sie sonst nicht sinnvoll referenziert werden können.

\subsection{Einreichung}

Beiträge werden elektronisch eingereicht.

Mit der Einreichung bestätigen die Autorinnen und Autoren, dass sie den Beitrag selbst verfasst haben, die Verantwortung für dessen Inhalt übernehmen und bereit sind, konstruktive Kritik als solche zu erkennen.

Die Redaktion bemüht sich um eine zügige Bearbeitung. Die tatsächliche Bearbeitungsdauer hängt jedoch von Faktoren ab, die gegenwärtig Gegenstand laufender Forschung sind.

Alle Korrespondenz erfolgt vorzugsweise schriftlich und in einem Tonfall, der die Wahrscheinlichkeit einer positiven Begutachtung nicht unnötig reduziert.

\section{Begutachtungsverfahren}

Alle Manuskripte werden von unabhängigen Fachgutachterinnen und Fachgutachtern geprüft.

Die Gutachten dienen der Sicherung wissenschaftlicher Qualität sowie der Unterhaltung der Redaktion.

Folgende Entscheidungen sind möglich:

\begin{itemize}
	\item Annahme
	\item Annahme nach Überarbeitung
	\item Überarbeitung mit erneuter Begutachtung
	\item Ablehnung
	\item ``Interessanter Gedanke, aber vielleicht lieber als Buch''
\end{itemize}

Die Entscheidung der Herausgeber ist endgültig, auch wenn sie gelegentlich überraschend wirkt.

\section{Veröffentlichungsrechte}

Mit Annahme des Beitrags räumen die Autorinnen und Autoren der Zeitschrift die erforderlichen Veröffentlichungsrechte ein.

Die Urheberrechte verbleiben grundsätzlich bei den Autorinnen und Autoren. Die Redaktion erhebt keinen Anspruch auf Eigentum an Ideen, Theorien oder innovativen Morgenroutinen.

\section{Kernaussagen}

Zu jedem Beitrag sollen drei bis fünf Kernaussagen eingereicht werden.

Diese sollen den wesentlichen Erkenntnisgewinn des Beitrags zusammenfassen und idealerweise verständlich sein, ohne dass die Leserin oder der Leser zuvor eine dreisemestrige Einführung in das Thema absolvieren muss.

\section{Korrekturabzüge}

Vor der Veröffentlichung erhalten die Autorinnen und Autoren einen Korrekturabzug.

Dieser dient der Korrektur von Fehlern und nicht der nachträglichen Neuerfindung des gesamten Beitrags.

Die Redaktion bittet darum, umfangreiche Grundsatzänderungen nach diesem Zeitpunkt auf zukünftige Forschungsprogramme zu verschieben.

\section{Anfragen der Autorinnen und Autoren}

Fragen zur Einreichung, Begutachtung oder Veröffentlichung können an die Redaktion gerichtet werden.

Bei Rückfragen zum Status eines Manuskripts empfiehlt sich ein Mindestabstand von sieben Tagen zwischen zwei Nachfragen. Die wiederholte Nachfrage im Abstand von 45 Minuten beschleunigt das Verfahren nach gegenwärtigem Kenntnisstand nicht.

\section*{Danksagungen}

Die Danksagungen bieten Raum für die Würdigung von Personen, Institutionen und Umständen, die zum Gelingen des Beitrags beigetragen haben.

Hierzu zählen insbesondere Kolleginnen und Kollegen, Diskussionspartner, Versuchspersonen, Bibliothekarinnen und Bibliothekare, Dr. Maximilian Methodius,, Koffeinquellen sowie alle weiteren Personen, deren Einfluss statistisch signifikant erscheint.

\section*{Interessenkonflikte}

Potenzielle Interessenkonflikte sind offenzulegen.

Hierzu zählen finanzielle Interessen, institutionelle Verpflichtungen, familiäre Bindungen oder langjährige Überzeugungen, die einer vollkommen objektiven Betrachtung möglicherweise im Wege stehen könnten.

Beispiel einer zulässigen Erklärung:

Die Autorinnen und Autoren erklären, dass keine Interessenkonflikte bestehen.

\section*{Anhänge (falls erforderlich)}

\subsection*{Anhang A}

Zusätzliche Materialien, Datensammlungen, Herleitungen oder andere Informationen, die für den Haupttext zu interessant, aber zu umfangreich sind.

\subsection*{Anhang B}

Weitere ergänzende Inhalte.

\bibliographystyle{zal}
\bibliography{ref}

\begin{biography}{beispielbild}
	Bitte geben Sie hier eine kurze Autorenbiografie mit maximal 120 Wörtern an.
	
	Die Biografie sollte Angaben zur fachlichen Tätigkeit, zu Forschungsschwerpunkten und gegebenenfalls zu bemerkenswerten Beiträgen zur angewandten Lebensführung enthalten.
	
	Ein Autorenfoto wird begrüßt. Bevorzugt werden Aufnahmen, die Kompetenz, Zuversicht und die grundsätzliche Fähigkeit zur erfolgreichen Bewältigung eines Montagvormittags vermitteln.
\end{biography}

\end{document}
